\chapter{2022年「大学物理1」\hdu 期末试题*10\%}
\date{2022年6月??日}{大学物理教学团队}{A0715011}
{未央物理讲师 \href{https://qm.qq.com/q/UPbGudx8cK}{\textbf{Axia}}}
{\qrcode{https://qm.qq.com/q/asJiytD5pC}}{「未央」HDU 物理营}
{\qrcode{http://weixin.qq.com/r/mhKYgFTEMQlOrRDj90eI}}{未央学社公众号}
\section{选择}
\begin{exercise}
两条金属轨道放在均匀磁场中.磁场方向垂直纸面向里，如图所示.在这两条轨道上垂直于轨道架设两条长而刚性的裸导线$P$与$Q$.金属线$P$中接入一个高阻伏特计.令导线$Q$保持不动，而导线$P$以恒定速度平行于导轨向左移动.(A)-(E)各图中哪一个正确表示伏特计电压$V$与时间上的关系？
\end{exercise}
\begin{solution}
恒定速度运动，感应电动势当然是常数啦，A，直线那个（图没画）
\end{solution}

\section{计算}
\vspace{1em}
\begin{minipage}{0.48\textwidth}
\begin{exercise}
如图所示，质量为$m$的木块，从高为$h$， 倾角为$θ$的光滑斜面上由静止开始下滑，滑入装着砂子的木箱中，砂子和木箱的总质量为$M$，木箱与一端固定, 劲度系数为$k$的水 平轻弹簧连接，最初弹簧为原长，木块落入后，弹簧的最大压缩量为 $l$，试求木箱与水平面间的摩擦系数$μ$.
\end{exercise}
\end{minipage}
\hfill
\begin{minipage}{0.48\textwidth}
\centering
\begin{tikzpicture}[scale=0.6]
\draw(1,3)--(5,3)--(1,5)--(1,3)--(0,3);
\draw(5,3)--(5,1)--(11,1)--(11,3);
\draw[very thick,紺桔梗](5.2,3)--(6.2,1)--(7.2,1)--(8.2,3);
\draw(7.7,2)--(8,2)--(8.3,2.3)--(8.5,1.5)--(9.3,2.3)--(9.5,1.5)--(10.3,2.3)--(10.5,1.5)--(11,2);
\draw[very thick,萌葱色](1,5)--(1.25,5.5)--(2.25,5)--(2,4.5)--(1,5);
\draw[<->](0.75,3)--(0.75,5);
\draw(0.5,5)--(1,5);
\coordinate[label=right:$m$] (m) at (2.25,5);
\coordinate[label=left:$h$] (h) at (0.75,4);
\coordinate[label=above:$M$] (M) at (6.5,3);
\coordinate[label=above:$k$] (k) at (9.5,2.5);
\end{tikzpicture}
\end{minipage}
\begin{solution}
\begin{enumerate}[(1)]
    \item 若不考虑物块撞击$M$的瞬间的竖直方向冲量对水平方向地面摩擦力冲量的影响（即$m$落入木箱时沿水平方向$m$与$M$间的冲力（内力）远大于地面与木箱的摩擦力（外力））\\
    以$M,\ m$为系统，水平方向动量守恒
    $$mv_0\cos\theta=\left(M+m\right)v$$
    由动能定理，$m$下滑过程中
    $$\frac{1}{2}mv_0^2=mgh$$
    解得
    $$v=\frac{m\sqrt{2gh}\cdot\cos\theta}{M+m}$$
    接下来的运动过程中，由能量守恒，物块和车的动能转化为和地面摩擦产生的热能和弹簧压缩时产生的弹性势能，即
    $$\frac{1}{2}\left(m+M\right)v^2=\mu\left(m+M\right)gl+\frac{1}{2}kl^2$$
    解得
    $$\mu=\frac{m^2h\cos^2\theta}{\left(M+m\right)^2l}-\frac{kl}{2\left(M+m\right)g}$$
    至此，本题可得全分.
    \item 题中并未强调\textbf{$m$落入木箱时沿水平方向$m$与$M$间的冲力（内力）远大于地面与木箱的摩擦力（外力）}，所以下面是对（1）情况下的修正.\footnote{感谢21级卓越学院吴同学提供的解答}\\
    对物块和车组成的系统，其水平方向和竖直方向上受力关系为
    $$F_{//}=\mu F_{\perp}$$
    记碰撞时间间隔为$\Delta t$，对上式子积分得
    $$\int_0^{\Delta t}F_{//}=\mu\int_0^{\Delta t} F_{\perp}$$
    即
    $$\Delta p_{//}=\mu\Delta p_{\perp}$$
    其中，地面对物块和车竖直方向上的冲量大小为$mv_0\cos\theta$（忽略重力）.
    对物块和车在水平方向上列动量定理
    $$mv_0\cos\theta-\mu mv_0\sin\theta=\left(m+M\right)v$$
    解得
    $$v=\frac{mv_0\cos\theta-\mu mv_0\sin\theta}{M+m}$$
    接下来到压缩最大过程中，由能量守恒
    $$\frac{1}{2}mv^2=\mu\left(M+m\right)gl+\frac{1}{2}kl^2$$
    最后解得
    $$\mu=\frac{\left(\frac{M+m}{m}\right)^2L+2h\sin\theta+\sqrt{\left[\left(\frac{M+m}{m}\right)^2L+2h\sin\theta \right]^2-\frac{4h\sin^2\theta}{g}\left(gh\cos^2\theta-\frac{M+m}{m^2}kl^2\right)}}{2h\sin^2\theta}$$
\end{enumerate}
\end{solution}